Without a full-rank hypothesis on the fixed point, the corner-restricted
fixed-point $*$-algebra of
Theorem~\ref{thm:cornerFixedPointsStarSubalgebra} still carries the block
structure through the compression onto the support sector; this is the
$\sum_kd_km_k\leq D$ support-sector form of the block representation in
Equation~(1.39) of~\cite{Wolf2012Quantum} invoked by
\cite[Theorem~6.14]{Wolf2012Quantum}. A corner-restricted fixed point
extended by zero on the complement of the support sector need not be a fixed
point of the ambient map, so the corner block form does not by itself give
the ambient fixed-point space a zero-block representation.

\begin{theorem}[Block form of the corner-restricted fixed points]%
    \label{thm:cornerFixedPoints_block_form}
    \lean{Kraus.cornerFixedPoints_blockDiagonal_iff}
    \leanok
    \uses{def:support_proj, def:kraus_map, def:kraus_adjoint_map,
        def:tp_kraus}
    Let $E(X)=\sum_iK_iXK_i^\dagger$ be a trace-preserving Kraus map on
    $\MN{D}$, with Heisenberg-picture adjoint
    $E^*(Y)=\sum_iK_i^\dagger YK_i$, let $\rho\succeq0$ satisfy
    $E(\rho)=\rho$, and let $Q$ be the support projection of $\rho$. There
    are a sector dimension $r\leq D$, a number $n$ of blocks, positive
    dimensions $d_0,\ldots,d_{n-1}$ and multiplicities
    $m_0,\ldots,m_{n-1}$ with $\sum_kd_km_k=r$, realized by an explicit
    identification of the index sets, and an isometry
    $W:\C^r\to\C^D$ with $W^\dagger W=\Id_r$ and $WW^\dagger=Q$, such that
    a matrix $Y\in\MN{D}$ satisfies $QYQ=Y$ and $QE^*(Y)Q=Y$ exactly when
    \begin{align}
        Y
         & =W\left(\bigoplus_{k=0}^{n-1}
        \Id_{m_k}\otimes B_k\right)W^\dagger
        \label{eq:spectral_ds_corner_element}
    \end{align}
    for some matrices $B_k\in\MN{d_k}$; that is, up to reordering the two
    tensor factors of each block,
    \begin{align}
        \{Y\in Q\MN{D}Q\mid QE^*(Y)Q=Y\}
         & =W\left(\bigoplus_{k=0}^{n-1}
        \Id_{m_k}\otimes\MN{d_k}\right)W^\dagger.
        \label{eq:spectral_ds_corner_blocks}
    \end{align}
    The right-hand side is the support-sector block representation in
    Equation~(1.39) of~\cite{Wolf2012Quantum}; the equivalence characterizes
    the corner-restricted set, not the ambient fixed-point space.
\end{theorem}

\begin{proof}
    \leanok
    \uses{thm:support_proj_fixed, thm:supported_kraus_corner_compression,
        thm:compression_on_support_posDef, thm:adjointFixedPoints_block_form}
    Since $\rho$ is a fixed point of $E$, the support projection satisfies
    $(\Id-Q)K_iQ=0$ by Theorem~\ref{thm:support_proj_fixed}, so the
    compressed family $QK_iQ$ is supported on the corner and
    trace-preserving there. Theorem~\ref{thm:supported_kraus_corner_compression}
    gives a support isometry $V:\C^r\to\C^D$ with
    $V^\dagger V=\Id_r$ and $VV^\dagger=Q$, together with a compressed family
    $C$ that is trace-preserving on $\C^r$. Moreover,
    $\sigma=V^\dagger\rho V$ is a positive definite fixed point of the
    compressed Schr\"odinger map by
    Theorem~\ref{thm:compression_on_support_posDef}. The block form of the
    Heisenberg-picture fixed points of $C$ from
    Theorem~\ref{thm:adjointFixedPoints_block_form} yields a unitary
    $U\in\MN{r}$, dimensions $d_k$, and multiplicities $m_k$ with
    $\sum_kd_km_k=r$.

    Write $\Phi(X):=VXV^\dagger$ for the compression isomorphism onto the
    corner and $C^*$ for the Heisenberg-picture adjoint of the compressed
    family. The corner-restricted fixed-point condition corresponds to the
    fixed-point equation of the compressed map:
    \begin{align}
        QE^*\!(\Phi(X))Q=\Phi(X)
         & \quad\Longleftrightarrow\quad C^*(X)=X.
        \notag
    \end{align}
    With $W:=VU$, which satisfies $W^\dagger W=\Id_r$ and
    $WW^\dagger=Q$, this equivalence gives
    (\ref{eq:spectral_ds_corner_element}) and
    (\ref{eq:spectral_ds_corner_blocks}). Finally, $r\leq D$ because the
    isometry $V$ embeds $\C^r$ into $\C^D$.
\end{proof}

\begin{theorem}[Weighted fixed points form a $*$-subalgebra]%
    \label{thm:weightedFixedPointsStarSubalgebra}
    \lean{Kraus.weightedFixedPointsStarSubalgebra}
    \leanok
    \uses{def:tp_kraus, thm:fixedPointsStarSubalgebra}
    Let $T(X)=\sum_iK_iXK_i^\dagger$ be trace-preserving, and let
    $\rho>0$ satisfy $T(\rho)=\rho$. If
    $F_T:=\{X\in\MN{D}\mid T(X)=X\}$, then
    $\rho^{-1/2}F_T\rho^{-1/2}$ is a $*$-subalgebra of $\MN{D}$.
\end{theorem}

\begin{proof}
    \leanok
    \uses{thm:fixedPointsStarSubalgebra}
    Set $B_i=\rho^{-1/2}K_i\rho^{1/2}$. The equation $T(\rho)=\rho$
    implies that the family $\{B_i\}$ is unital, and trace preservation of
    $T$ implies that $\rho$ is a positive-definite fixed point of the
    adjoint map of the gauged family. Apply
    Theorem~\ref{thm:fixedPointsStarSubalgebra} to $B$. The equivalence
    \begin{align}
        \sum_iB_iXB_i^\dagger=X
         & \quad\Longleftrightarrow\quad
        T(\rho^{1/2}X\rho^{1/2})=\rho^{1/2}X\rho^{1/2}
        \notag
    \end{align}
    identifies the fixed points of the gauged map with
    $\rho^{-1/2}F_T\rho^{-1/2}$.
\end{proof}

For a singular fixed point the conjugation is taken with the inverse square
root on the support of $\rho$, and the conjugated set lives in the corner
algebra of the support projection. The corollary in~\cite{Wolf2012Quantum}
takes a maximum-rank fixed point, for which every fixed point of $T$ is
supported on the support of $\rho$; the following two statements take an
arbitrary positive semidefinite fixed point and restrict to the fixed points
supported on the support of $\rho$.
Theorem~\ref{thm:maximalSupport_fixedPoint} below supplies a fixed point
whose support carries every fixed point, and at such a fixed point the
restriction disappears
(Theorem~\ref{thm:maximalSupport_weightedCorner}).
% Maintainer note: the restriction history is recorded in
% docs/paper-gaps/wolf_cor67_maximal_support_restriction.tex.

\begin{theorem}[Weighted corner fixed points form a $*$-subalgebra]%
    \label{thm:weightedCornerFixedPointsStarSubalgebra}
    \lean{Kraus.weightedCornerFixedPointsStarSubalgebra}
    \leanok
    \uses{def:tp_kraus, def:support_proj, rem:corner_algebra_star_instances}
    Let $T(X)=\sum_iK_iXK_i^\dagger$ be trace-preserving, let
    $\rho\succeq0$ satisfy $T(\rho)=\rho$, and let $Q$ be the support
    projection of $\rho$. The corner elements $Y\in Q\MN{D}Q$ such that
    $\sqrt{\rho}\,Y\sqrt{\rho}$ is a fixed point of $T$ form a
    $*$-subalgebra of the corner algebra $Q\MN{D}Q$.
\end{theorem}

\begin{proof}
    \leanok
    \uses{thm:weightedFixedPointsStarSubalgebra,
        thm:supported_kraus_corner_compression,
        thm:compression_on_support_posDef, thm:support_proj_fixed}
    Closure under addition, scalars, and conjugation is linearity together
    with the stability of fixed points under conjugate transposition. The
    corner unit $Q$ belongs to the set because
    $\sqrt{\rho}\,Q\sqrt{\rho}=\rho$ is a fixed point. The support projection
    absorbs the square root, $Q\sqrt{\rho}=\sqrt{\rho}=\sqrt{\rho}\,Q$,
    since
    \begin{align}
        (\Id-Q)\sqrt{\rho}
        ((\Id-Q)\sqrt{\rho})^\dagger
         & =(\Id-Q)\rho(\Id-Q)=0.
        \notag
    \end{align}

    For closure under multiplication, compress to the support sector. Since
    $\rho$ is a fixed point of $T$, the support projection satisfies
    $(\Id-Q)K_iQ=0$ by Theorem~\ref{thm:support_proj_fixed}. The compressed
    family $QK_iQ$ is trace-preserving on the corner.
    Theorem~\ref{thm:supported_kraus_corner_compression} gives a support
    isometry $V:\C^r\to\C^D$ with $V^\dagger V=\Id_r$ and
    $VV^\dagger=Q$, together with a trace-preserving compressed
    family $C$ on $\C^r$ whose Schr\"odinger map has the positive definite
    fixed point $\sigma=V^\dagger\rho V$ by
    Theorem~\ref{thm:compression_on_support_posDef}. The square root
    compresses along the isometry:
    \begin{align}
        \sqrt{\sigma}
         & =V^\dagger\sqrt{\rho}\,V,
        \label{eq:spectral_ds_sqrt_compression}
    \end{align}
    by uniqueness of the positive square root, since
    $(V^\dagger\sqrt{\rho}V)^2
        =V^\dagger\sqrt{\rho}\,Q\sqrt{\rho}V=\sigma$.
    Hence, for corner-supported $Y$,
    \begin{align}
        \sqrt{\sigma}\,(V^\dagger YV)\sqrt{\sigma}
         & =V^\dagger(\sqrt{\rho}\,Y\sqrt{\rho})V.
        \label{eq:spectral_ds_weighted_transport}
    \end{align}
    For corner-supported $W$, the compressed map satisfies
    \begin{align}
        C(V^\dagger WV)
         & =V^\dagger QT(W)QV
        =V^\dagger T(W)V,
        \notag
    \end{align}
    since $T(W)$ is again corner-supported. Thus $T(W)=W$ exactly when the
    compressed map fixes $V^\dagger WV$. The set therefore corresponds,
    element by element, to the weighted fixed points of the compressed
    family with respect to the positive definite fixed point $\sigma$,
    which form a $*$-subalgebra by
    Theorem~\ref{thm:weightedFixedPointsStarSubalgebra}.

    Closure under multiplication transports back through the correspondence.
    A corner-supported $Y_j$ satisfies $Y_jQ=Y_j$, so
    \begin{align}
        (V^\dagger Y_1V)(V^\dagger Y_2V)
         & =V^\dagger Y_1(VV^\dagger)Y_2V
        =V^\dagger(Y_1Q)Y_2V
        =V^\dagger(Y_1Y_2)V.
        \notag
    \end{align}
    The product of two members therefore corresponds to the product of
    their compressed images, which lies in the compressed
    $*$-subalgebra.
\end{proof}

Combined with the conjugation correspondence of
Theorem~\ref{thm:weightedCorner_sqrt_surjective}, the carrier of this
$*$-subalgebra realizes
\begin{align}
    \rho^{-1/2}
    \{X\in\MN{D}\mid T(X)=X,\ QXQ=X\}
    \rho^{-1/2},
    \notag
\end{align}
with the inverse square root taken on the support of $\rho$: the singular
case of Corollary~6.7 of~\cite{Wolf2012Quantum}, restricted to the fixed
points supported on the support of $\rho$.

\begin{theorem}[Conjugation by the square root is onto the corner-supported fixed
        points]%
    \label{thm:weightedCorner_sqrt_surjective}
    \lean{Kraus.exists_weightedCorner_sqrt_eq_of_fixedPoint}
    \leanok
    \uses{def:tp_kraus, def:support_proj}
    In the setting of
    Theorem~\ref{thm:weightedCornerFixedPointsStarSubalgebra}, every fixed
    point $X$ of $T$ with $QXQ=X$ arises as
    $X=\sqrt{\rho}\,Y\sqrt{\rho}$ for a corner-supported $Y$ with
    $\sqrt{\rho}\,Y\sqrt{\rho}$ fixed by $T$.
\end{theorem}

\begin{proof}
    \leanok
    \uses{thm:supported_kraus_corner_compression,
        thm:compression_on_support_posDef, thm:support_proj_fixed}
    Use Theorem~\ref{thm:supported_kraus_corner_compression} to compress to
    the support sector as in the proof of
    Theorem~\ref{thm:weightedCornerFixedPointsStarSubalgebra}, and set
    \begin{align}
        Y
         & =V\left(\sqrt{\sigma}^{-1}(V^\dagger XV)
        \sqrt{\sigma}^{-1}\right)V^\dagger,
        \label{eq:spectral_ds_weighted_preimage}
    \end{align}
    where $\sigma=V^\dagger\rho V$ is positive definite, so
    $\sqrt{\sigma}$ is invertible. Then $Y$ is corner-supported. Writing
    $Z=\sqrt{\sigma}^{-1}(V^\dagger XV)\sqrt{\sigma}^{-1}$ for the
    compressed middle factor, (\ref{eq:spectral_ds_weighted_transport})
    gives
    \begin{align}
        V^\dagger(\sqrt{\rho}\,Y\sqrt{\rho})V
         & =\sqrt{\sigma}\,(V^\dagger YV)\sqrt{\sigma}
        =\sqrt{\sigma}\,Z\sqrt{\sigma}
        =V^\dagger XV.
        \notag
    \end{align}
    Hence $\sqrt{\rho}\,Y\sqrt{\rho}=QXQ=X$, since both sides are
    corner-supported. In particular, $\sqrt{\rho}\,Y\sqrt{\rho}$ is fixed
    by $T$.
\end{proof}

The support hypothesis on the fixed points is discharged by exhibiting a
fixed point whose support carries every fixed point. The first ingredient is
that Loewner domination transfers the support.

\begin{lemma}[The transfer map of a trace-preserving family is a channel]%
    \label{lem:isChannel_transferMap}
    \lean{Kraus.isChannel_transferMap}
    \leanok
    \uses{def:channel, def:tp_kraus}
    Let $K_0,\ldots,K_{d-1}$ be a trace-preserving Kraus family. The map
    $E(X)=\sum_iK_iXK_i^\dagger$ is a quantum channel.
\end{lemma}

\begin{proof}\leanok
    The transfer-map formula is the finite Kraus map. Its Kraus form is
    completely positive, while the normalization
    $\sum_iK_i^\dagger K_i=\Id$ gives $\tr(E(X))=\tr(X)$. Hence $E$ is a
    channel.
\end{proof}

\begin{lemma}[Scalar Loewner domination transfers the support]%
    \label{lem:stationaryProj_absorb_of_le}
    \lean{Kraus.stationaryProj_absorb_of_le_smul}
    \leanok
    \uses{def:support_proj}
    Let $\rho,P\succeq0$ and $c\in\C$ with $P\preceq c\rho$, and let $Q$
    be the support projection of $\rho$. Then $QP=P$ and $PQ=P$.
\end{lemma}

\begin{proof}
    \leanok
    \uses{def:support_proj}
    The support projection absorbs $\rho$, so $(\Id-Q)\rho=0$ and
    \begin{align}
        0
         & \preceq(\Id-Q)P(\Id-Q)
        \preceq(\Id-Q)(c\rho)(\Id-Q)
        =0,
        \notag
    \end{align}
    because conjugation by $\Id-Q$ preserves the Loewner order. Hence
    \begin{align}
        ((\Id-Q)\sqrt{P})
        ((\Id-Q)\sqrt{P})^\dagger
         & =(\Id-Q)P(\Id-Q)=0,
        \notag
    \end{align}
    so $(\Id-Q)\sqrt{P}=0$ and
    $(\Id-Q)P=((\Id-Q)\sqrt{P})\sqrt{P}=0$. Thus $QP=P$; taking conjugate
    transposes gives $PQ=P$ since $P$ and $Q$ are Hermitian.
\end{proof}

\begin{theorem}[A fixed point of maximal support from bounded orbits]%
    \label{thm:maximalSupport_fixedPoint_of_bounded_orbits}
    \lean{IsPositiveMap.exists_maximalSupport_fixedPoint_of_hasBoundedOrbits}
    \leanok
    \uses{def:positive_map, def:has_bounded_orbits,
        def:mean_ergodic_projection, def:support_proj}
    Let $T:\MN{D}\to\MN{D}$ be positive and have bounded orbits, and let
    $T_\infty$ be its mean-ergodic projection. Then
    $\rho_0=T_\infty(\Id)$ is positive semidefinite and fixed by $T$. If
    $Q_0$ is the support projection of $\rho_0$, then every fixed point $X$
    of $T$ satisfies $Q_0XQ_0=X$. This is the bounded-orbit form of the
    maximal-support argument underlying
    \cite[Proposition~6.9]{Wolf2012Quantum}.
\end{theorem}

\begin{proof}
    \leanok
    \uses{thm:positive_mean_ergodic_projection,
        thm:mean_ergodic_bounded_orbits,
        thm:positive_map_conjTranspose, thm:psd_le_trace_smul_one,
        lem:stationaryProj_absorb_of_le}
    Positivity of the mean-ergodic projection gives $\rho_0\succeq0$, while
    its range and idempotence give $T(\rho_0)=\rho_0$. If $A\succeq0$, then
    $A\preceq\tr(A)\Id$; applying the positive map $T_\infty$ gives
    \begin{align}
        \tr(A)\rho_0-T_\infty(A)
         & =T_\infty\!\left(\tr(A)\Id-A\right)\succeq0,
         &
        T_\infty(A)
         & \preceq\tr(A)\rho_0.
        \notag
    \end{align}
    Scalar domination therefore gives
    $Q_0T_\infty(A)Q_0=T_\infty(A)$. For a fixed Hermitian matrix $H$,
    write $H=H^+-H^-$. Since $T_\infty(H)=H$, linearity gives
    \begin{align}
        H
         & =T_\infty(H^+)-T_\infty(H^-),
        \notag
    \end{align}
    so $Q_0HQ_0=H$. Finally, decompose an arbitrary fixed point into two
    Hermitian fixed matrices, using preservation of conjugate transpose, and
    conclude that $Q_0XQ_0=X$.
\end{proof}

\begin{theorem}[A fixed point of maximal support]%
    \label{thm:maximalSupport_fixedPoint}
    \lean{IsPositiveMap.exists_maximalSupport_fixedPoint}
    \leanok
    \uses{def:positive_map, def:trace_preserving, def:mean_ergodic_projection,
        def:support_proj}
    Let $T:\MN{D}\to\MN{D}$ be positive and trace-preserving, and let
    $T_\infty$ be its mean-ergodic projection. Then
    $\rho_0=T_\infty(\Id)$ is positive semidefinite and fixed by $T$. With
    $Q_0$ the support projection of $\rho_0$, every fixed point $X$ of $T$
    satisfies $Q_0XQ_0=X$. This is the maximal-support property of fixed
    points in Section~6.4 of~\cite{Wolf2012Quantum}.
\end{theorem}

\begin{proof}
    \leanok
    \uses{thm:positive_trace_preserving_mean_ergodic_projection,
        thm:maximalSupport_fixedPoint_of_bounded_orbits}
    Theorem~\ref{thm:positive_trace_preserving_mean_ergodic_projection}
    gives bounded orbits. Apply
    Theorem~\ref{thm:maximalSupport_fixedPoint_of_bounded_orbits}.
\end{proof}

\begin{theorem}[Restriction to the support of a stationary positive matrix]%
    \label{thm:stationarySupport_restriction}
    \lean{IsPositiveMap.map_supported_on_fixedPoint_support}
    \leanok
    \uses{def:positive_map, def:support_proj,
        lem:stationaryProj_absorb_of_le}
    Let $T:\MN{D}\to\MN{D}$ be positive, and let
    $\rho\succeq0$ satisfy $T(\rho)=\rho$. If $Q$ is the support projection
    of $\rho$, then, for every $X\in\MN{D}$,
    \begin{align}
        QXQ=X
         & \quad\Longrightarrow\quad QT(X)Q=T(X).
        \label{eq:spectral_ds_support_invariance}
    \end{align}
\end{theorem}

\begin{proof}
    \leanok
    \uses{lem:stationaryProj_absorb_of_le}
    First suppose $A\succeq0$ and $QAQ=A$. On the support of $\rho$, the
    matrix $\rho$ is positive definite, so $Q\preceq c\rho$ for some
    positive scalar $c$. Moreover, $A\preceq\tr(A)Q$, because this follows
    by compressing $\tr(A)\Id-A\succeq0$ with $Q$. Positivity and
    $T(\rho)=\rho$ therefore give
    \begin{align}
        0
         & \preceq T(A)\preceq c\tr(A)\rho.
        \notag
    \end{align}
    Scalar domination transfers the support, hence $QT(A)Q=T(A)$.

    Write $H_1=X+X^*$ and $H_2=\mathrm{i}(X-X^*)$, so that both matrices
    are Hermitian and $X=\tfrac12H_1-\tfrac{\mathrm{i}}2H_2$. Setting
    $A_j^\pm=QH_j^\pm Q$ gives positive matrices with
    $QA_j^\pm Q=A_j^\pm$ and
    \begin{align}
        X
         & =\tfrac12(A_1^+-A_1^-)
        -\tfrac{\mathrm{i}}2(A_2^+-A_2^-).
        \notag
    \end{align}
    Applying the positive-input conclusion to these four matrices and using
    linearity of $T$ gives (\ref{eq:spectral_ds_support_invariance}).
\end{proof}

\begin{theorem}[Restriction to full-rank fixed points]%
    \label{thm:stationarySupport_compression}
    \lean{IsPositiveMap.exists_maximalSupportCompression}
    \leanok
    \uses{thm:maximalSupport_fixedPoint, thm:stationarySupport_restriction}
    Let $T:\MN{D}\to\MN{D}$ be positive and trace preserving, set
    $\rho_0=T_\infty(\Id)$, and let $Q_0$ be the support projection of
    $\rho_0$. There are a finite-dimensional Hilbert space $\mathcal H$ and
    an isometry $V:\mathcal H\to\C^D$ with
    $V^*V=\Id_{\mathcal H}$ and $VV^*=Q_0$ such that
    \begin{align}
        \widetilde T(Y)
         & =V^*T(VYV^*)V
        \label{eq:spectral_ds_support_compression}
    \end{align}
    is positive and trace preserving and has a positive definite fixed
    point. Moreover,
    \begin{align}
        T(VYV^*)
         & =V\widetilde T(Y)V^*,
        \label{eq:spectral_ds_compression_intertwining} \\
        T(X)=X
         & \quad\Longleftrightarrow\quad
        \exists\,Y,\quad
        \widetilde T(Y)=Y\ \text{ and }\ X=VYV^*.
        \label{eq:spectral_ds_compression_fixed}
    \end{align}
    These are Equations~(6.52) and~(6.51), respectively,
    of~\cite{Wolf2012Quantum}. Thus the complementary fixed-point summand
    vanishes.
\end{theorem}

\begin{proof}
    \leanok
    \uses{thm:maximalSupport_fixedPoint, thm:stationarySupport_restriction}
    Theorem~\ref{thm:maximalSupport_fixedPoint} gives $\rho_0$ and shows
    that $Q_0XQ_0=X$ for every fixed point $X$ of $T$. Choose an isometry
    onto the range of $Q_0$. Positivity follows from positivity of
    conjugation and of $T$. The matrix $VYV^*$ is supported on $Q_0$, so
    Theorem~\ref{thm:stationarySupport_restriction} shows that its image is
    also supported there. Inserting $Q_0=VV^*$ gives
    \begin{align}
        T(VYV^*)
         & =Q_0T(VYV^*)Q_0
        =VV^*T(VYV^*)VV^*
        =V\widetilde T(Y)V^*,
        \notag
    \end{align}
    which proves (\ref{eq:spectral_ds_compression_intertwining}). By
    cyclicity of the trace, the isometry identity, and trace preservation of
    $T$,
    \begin{align}
        \tr(\widetilde T(Y))
         & =\tr(V\widetilde T(Y)V^*)
        =\tr(T(VYV^*))
        =\tr(VYV^*)
        =\tr(Y).
        \notag
    \end{align}
    Thus $\widetilde T$ is trace preserving. The compressed point
    $\sigma_0=V^*\rho_0V$ is positive definite. Since
    $V\sigma_0V^*=Q_0\rho_0Q_0=\rho_0$,
    \begin{align}
        \widetilde T(\sigma_0)
         & =V^*T(\rho_0)V
        =V^*\rho_0V
        =\sigma_0.
        \notag
    \end{align}
    If $T(X)=X$, then $Q_0XQ_0=X$, and hence
    $X=V(V^*XV)V^*$. Moreover,
    \begin{align}
        \widetilde T(V^*XV)
         & =V^*T(V(V^*XV)V^*)V
        =V^*T(X)V
        =V^*XV.
        \notag
    \end{align}
    Conversely, if $\widetilde T(Y)=Y$, then
    (\ref{eq:spectral_ds_compression_intertwining}) gives
    $T(VYV^*)=V\widetilde T(Y)V^*=VYV^*$. These two implications prove
    (\ref{eq:spectral_ds_compression_fixed}).
\end{proof}

The next three theorems are generic finite-matrix facts. They do not depend
on positive maps, fixed points, or Theorem~6.14 of~\cite{Wolf2012Quantum}.
